% =====================================================================
%  CARNET D'ENTRAÎNEMENT — CHIMIE — FICHE 3 — THÈME 2
%  Composition : Z, A, notation, atomes et ions — TUTO
% =====================================================================
\documentclass[11pt,a4paper]{article}

\usepackage[utf8]{inputenc}
\usepackage[T1]{fontenc}
\usepackage{lmodern}
\usepackage[french]{babel}

\usepackage{amsmath,amssymb,mathtools}
\usepackage{xcolor}
\usepackage{colortbl}
\usepackage{tikz}
\usetikzlibrary{arrows.meta,positioning,calc}
\usepackage[most]{tcolorbox}
\usepackage{enumitem}
\usepackage{booktabs}
\usepackage{array}
\usepackage[a4paper,margin=2.1cm,headheight=26pt,headsep=9pt]{geometry}
\usepackage{fancyhdr}
\usepackage{titlesec}
\titlespacing*{\section}{0pt}{6pt}{3pt}
\titleformat{\section}{\large\bfseries\color{primarydark}}{}{0pt}{}
\usepackage[hidelinks]{hyperref}

\definecolor{primary}{RGB}{124,78,140}
\definecolor{primarydark}{RGB}{74,44,92}
\definecolor{defcol}{RGB}{0,114,178}
\definecolor{thmcol}{RGB}{0,140,120}
\definecolor{formulacol}{RGB}{198,93,9}
\definecolor{attncol}{RGB}{200,40,55}
\definecolor{protoncol}{RGB}{205,50,50}
\definecolor{neutroncol}{RGB}{105,108,120}
\definecolor{electroncol}{RGB}{25,95,205}

\newcommand{\nuc}[3]{\prescript{#1}{#2}{\mathrm{#3}}}
\newcommand{\pp}{p^{+}}
\newcommand{\ee}{e^{-}}
\newcommand{\nz}{n^{0}}
\newcommand{\np}{n(\pp)}
\newcommand{\nnn}{n(\nz)}

\pagestyle{fancy}
\fancyhf{}
\fancyhead[L]{\small\textcolor{primarydark}{\textbf{Fiche 3 — Thème 2}} : Composition des atomes et des ions}
\fancyhead[R]{\small\textcolor{primarydark}{\textsc{Tuto}}}
\fancyfoot[C]{\small\thepage}
\renewcommand{\headrulewidth}{0.8pt}
\renewcommand{\headrule}{\hbox to\headwidth{\color{primary}\leaders\hrule height \headrulewidth\hfill}}

\tcbset{
  boxbase/.style={enhanced,breakable,boxrule=0.9pt,arc=2.5pt,
     left=3mm,right=3mm,top=2.2mm,bottom=2.2mm,
     fonttitle=\bfseries,coltitle=white,toptitle=0.5mm,bottomtitle=0.5mm},
}
\newtcolorbox{cle}[1][]{boxbase,colback=thmcol!6,colframe=thmcol,
  title={\textsc{Point clé}\ifx\relax#1\relax\relax\else\textmd{ : \itshape #1}\fi}}
\newtcolorbox{formule}[1][]{boxbase,colback=formulacol!7,colframe=formulacol,
  title={\textsc{Formule}\ifx\relax#1\relax\relax\else\textmd{ : \itshape #1}\fi}}
\newtcolorbox{piege}[1][]{boxbase,colback=attncol!7,colframe=attncol,
  title={\textsc{Piège à éviter}\ifx\relax#1\relax\relax\else\textmd{ : \itshape #1}\fi}}
\newtcolorbox{meth}[1][]{boxbase,colback=defcol!5,colframe=defcol,
  title={\textsc{Méthode}\ifx\relax#1\relax\relax\else\textmd{ : \itshape #1}\fi}}

\setlist{leftmargin=1.5em,topsep=2pt,itemsep=1.5pt}

\begin{document}

\begin{tikzpicture}
\node[rectangle,rounded corners=6pt,fill=primary,text=white,
      minimum width=\textwidth,minimum height=1.5cm,align=center,inner sep=6pt]
  {\Large\bfseries Thème 2 — Composition des atomes et des ions\\[2pt]
   \normalsize\mdseries numéro atomique $Z$, nombre de masse $A$, comptage de $\pp$, $\nz$, $\ee$};
\end{tikzpicture}

\vspace{3pt}
{\small\itshape Objectif : à partir de la notation $\nuc{A}{Z}{X}$ (et d'une éventuelle charge), compter
sans erreur le nombre de \textbf{protons}, de \textbf{neutrons} et d'\textbf{électrons} d'un atome ou d'un
ion.}

\section*{1. Deux nombres suffisent : $Z$ et $A$}

\begin{formule}[les trois relations fondamentales]
\[ \boxed{\,Z = \np\,} \qquad \boxed{\,A = \np + \nnn\,} \qquad \boxed{\,\nnn = A - Z\,} \]
\begin{itemize}
  \item $Z$ = \textbf{numéro atomique} = nombre de protons. Il définit \emph{l'identité} de l'élément
        (tout carbone a $Z=6$, tout oxygène $Z=8$\dots).
  \item $A$ = \textbf{nombre de masse} = nombre total de \textbf{nucléons} (protons + neutrons).
  \item Le nombre de neutrons s'obtient par \textbf{différence} : $\nnn = A - Z$. On a toujours $A \ge Z$.
\end{itemize}
\end{formule}

\begin{center}
\begin{tikzpicture}[scale=1.0]
  \node[font=\Huge] (X) at (0,0) {$\mathrm{X}$};
  \node[font=\large,protoncol,anchor=east] (A) at (-0.25,0.42) {$A$};
  \node[font=\large,defcol,anchor=east]   (Z) at (-0.25,-0.42) {$Z$};
  \draw[<-,>=Stealth,protoncol] (A.north west) .. controls (-1.6,1.2) .. (-3.1,0.9)
       node[left,font=\footnotesize,align=right,protoncol]{nombre de masse $A$\\ \footnotesize (en haut)};
  \draw[<-,>=Stealth,defcol] (Z.south west) .. controls (-1.6,-1.2) .. (-3.1,-0.9)
       node[left,font=\footnotesize,align=right,defcol]{numéro atomique $Z$\\ \footnotesize (en bas)};
  \draw[<-,>=Stealth,black!70] (X.east) .. controls (1.4,0.0) .. (2.8,0.0)
       node[right,font=\footnotesize,align=left]{symbole de\\ l'élément};
  \node[font=\footnotesize,align=center] at (0,-1.7)
       {Exemple : $\nuc{35}{17}{Cl}$ $\Rightarrow$ $17$ protons, $35-17=18$ neutrons.};
\end{tikzpicture}
\end{center}

\section*{2. De l'atome neutre à l'ion}

Un atome \textbf{neutre} a autant d'électrons que de protons : $\ee = Z$. Un \textbf{ion} a gagné ou perdu
des électrons ; \textbf{seuls les électrons changent}.

\begin{formule}[nombre d'électrons]
\[ \text{atome neutre : } \boxed{\,\ee = Z\,} \qquad\qquad
   \text{ion de charge } q : \boxed{\,\ee = Z - q\,} \]
$q$ est \emph{signé} : $q=+2$ pour $\mathrm{Mg}^{2+}$ (on \emph{retire} $2$ électrons),
$q=-2$ pour $\mathrm{O}^{2-}$ (on \emph{ajoute} $2$ électrons). Un \textbf{cation} ($+$) a perdu des
électrons, un \textbf{anion} ($-$) en a gagné.
\end{formule}

\begin{cle}[le noyau est intouchable]
Quand un atome devient un ion, $Z$ (protons) et $\nnn$ (neutrons) \textbf{ne changent pas} : le noyau
est identique. $\mathrm{Ca}$ et $\mathrm{Ca}^{2+}$ ont le même noyau ($20\,\pp$, $20\,\nz$) et ne
diffèrent que par les électrons ($20$ contre $18$).
\end{cle}

\begin{meth}[compter $\pp$, $\nz$, $\ee$ d'une espèce $\nuc{A}{Z}{X}^{q}$]
\textbf{1.} Protons : $\np = Z$ (toujours). \quad
\textbf{2.} Neutrons : $\nnn = A - Z$ (toujours ; le noyau ne bouge pas). \quad
\textbf{3.} Électrons : $\ee = Z - q$ ($q$ signé).
\end{meth}

\begin{piege}
\textbf{« Combien d'électrons dans le noyau ? »} $\to$ toujours \textbf{0} : le noyau ne contient que des
nucléons (protons + neutrons). \; \textbf{Ne jamais modifier $A$ ni $Z$ pour un ion} : la charge change
seulement le nombre d'électrons. \; \textbf{$A$ n'est pas le nombre d'électrons} : c'est le nombre de
\emph{nucléons}.
\end{piege}

\end{document}
